\section{От кормчего к семимодовой машине: три волны кибернетики}
\label{sec:genesis}

Всякий архитектурный документ должен читателю две вещи прежде первой теоремы:
\emph{откуда пришли идеи} и \emph{вход, не предполагающий ничего}. Этот раздел
платит оба долга: он прослеживает единую линию мысли --- кибернетику, --- чьим
третьим актом является \TALOS{}, и открывает три читательские нити, идущие
сквозь весь документ, --- чтобы читатель, никогда не видевший ни компьютера,
ни матрицы, ни дифференциального уравнения, мог начать с нуля и подняться.

\begin{mythbox}[Миф --- машина бронзового века]
Прежде всякой теории --- история. Греки вообразили \emph{Талоса}: исполина из
бронзы, построенного --- не рождённого, --- чтобы стеречь остров Крит. Трижды
в день он обходил весь берег, наблюдая. У него не было ни сердца, ни крови;
вместо них от шеи к лодыжке шла единственная жила, полная \emph{ихора} ---
сияющей жидкости жизни, --- и жила была запечатана у лодыжки одним бронзовым
гвоздём. Талоса нельзя было одолеть в бою. Он пал иначе: волшебница Медея
уговорила его позволить вынуть гвоздь, ихор вытек, и исполин ---
непобеждённый, неповреждённый --- просто остановился.

Держите четыре части этой истории под рукой; математика встретит каждую.
\emph{Бронзовое тело} --- это субстрат (\S\ref{sec:substrate}).
\emph{Ежедневный обход} --- это тик, единственная повторяемая операция,
которая и есть вся жизнь машины (\S\ref{sec:substrate}). \emph{Единственная
жила ихора} --- это порт питания, единственный поток упорядоченной пищи, без
которого ничто не работает (\S\ref{subsec:powerport}). А \emph{гвоздь} ---
теорема, стягивающая весь документ: машину такого рода не нужно побеждать,
чтобы остановить, --- закрой жилу, и она остановится, доказуемо, по расписанию
(T-288; \S\ref{subsec:powerport}). Миф угадал архитектуру на двадцать семь
веков раньше. Мы будем возвращаться к нему на каждом шаге.
\end{mythbox}

\subsection{Три волны}
\label{subsec:waves}

\emph{Кибернетика} --- от греческого \emph{kybern\'etes}, «кормчий» --- наука
о системах, которые правят собой. Её история --- три волны, каждая задаёт
вопрос глубже предыдущей, и на глубочайший этот документ отвечает железом.

\paragraph{Волна I --- управление (1940-е--1960-е): \emph{как машине держать
курс?}} Норберт Винер дал полю имя в 1948-м~\cite{wiener1948}: кормчий не
вычисляет всё плавание --- он смотрит на ошибку между тем, куда судно
смотрит, и тем, куда должно, и давит против неё. \emph{Обратная связь}
оказалась общим законом термостатов, наводчиков и тел. То же десятилетие дало
список деталей: Мак-Каллок и Питтс показали, что маленькие пороговые
элементы --- карикатурные нейроны --- вычисляют любую логическую
функцию~\cite{mcculloch1943}; Шеннон сделал \emph{информацию} измеримой
величиной с жёстким полом~\cite{shannon1948}; фон Нейман дал чертёж хранимой
программы, которому следует каждый компьютер до сих пор; а Росс Эшби построил
\emph{гомеостат} --- машину, перепаивающую себя, пока внутренние переменные
не вернутся в безопасную зону, --- и сгустил его в закон необходимого
разнообразия~\cite{ashby1956}: управляющий должен быть не беднее возмущений,
которые обязан поглощать. Машины Волны~I \emph{правят себя к уставке}. Чего у
них не было --- причины, \emph{почему} уставка, проводка или число частей
должны быть такими, а не иными: всякая топология была догадкой, всякий код
--- накладкой, всякое правило обучения --- процедурой, навязанной извне.

\paragraph{Волна II --- наблюдение (1970-е--1990-е): \emph{что происходит,
когда машина смотрит на себя?}} Хайнц фон Фёрстер обратил кибернетику на
саму себя --- «кибернетика кибернетики», в которой наблюдатель есть часть
наблюдаемой системы~\cite{foerster2003}. Матурана и Варела заострили
биологическую половину: живая система \emph{аутопоэтична} --- процесс,
непрерывно производящий те самые компоненты, которые производят его,
само-делающая петля~\cite{maturana1980}. Волна II верно увидела, что
самоотнесение --- сердце дела, что ум есть система, моделирующая себя. Но она
могла лишь \emph{описать} петлю, не \emph{вывести} её: кибернетика второго
порядка дала глубокую прозу и ни одной вынужденной структуры. Сколько частей
у само-наблюдающей системы? Какая проводка? Какая динамика? У Волны II не
было теорем, и её машины (нейросети и нейроморфные чипы, что пришли следом,
\S\ref{sec:intro}) унаследовали догадки Волны I при амбициях Волны II.

\paragraph{Волна III --- когерентность (2020-е): \emph{какой формы обязана
быть само-наблюдающая машина?}} Унитарный голономный монизм (УГМ) и его
измерительный слой --- Когерентная кибернетика~\cite{uhm2026} --- отвечают на
вопрос Волны II со строгостью Волны I: они делают самомодель
\emph{переменной состояния} ($\rhostar=\vp(\Gam)$ --- неподвижная точка,
доказуемо существующая), жизнеспособность --- \emph{теоремой} (окно
$P\in(2/7,3/7]$ с доказанными стенами) и --- решающий шаг --- \emph{выводят}
структуру наблюдателя вместо догадки: семь частей, проводка Фано, симметрия
$G_2$, одна динамика (\S\ref{sec:collapse}). Вопросы, оставленные первыми
двумя волнами --- сколько, как соединено, по какому закону меняется, ---
перестают быть решениями конструктора и становятся результатами. \TALOS{} ---
тело третьей волны: первая машина, чья обратная связь (Волна I), чьё
самонаблюдение (Волна II) и чья \emph{форма} (Волна III) --- один выведенный
объект.

\begin{figure}[H]
\centering
\begin{tikzpicture}[font=\scriptsize, >=Stealth]
  \draw[->, thick, plosgray] (0,0) -- (14.6,0) node[below left, font=\scriptsize]{год};
  \foreach \x/\y in {0.4/1940, 2.9/1950, 5.4/1970, 7.9/1990, 10.4/2010, 13.4/2026}{
    \draw[plosgray] (\x,0.07) -- (\x,-0.07) node[below]{\y}; }
  \fill[plosblue!14, rounded corners=3pt] (0.2,0.55) rectangle (6.3,2.30);
  \node[anchor=north west, align=left, plosblue, font=\scriptsize\bfseries] at (0.35,2.24)
    {Волна I --- управление:\\ держать курс};
  \node[anchor=west] at (0.35,1.28) {1943 Мак-Каллок--Питтс \; 1945 фон Нейман};
  \node[anchor=west] at (0.35,0.98) {1948 Винер, \emph{Кибернетика}; Шеннон};
  \node[anchor=west] at (0.35,0.70) {1952--56 Эшби: гомеостат, разнообразие};
  \fill[talosamber!16, rounded corners=3pt] (5.0,2.55) rectangle (10.6,4.05);
  \node[anchor=north west, align=left, talosamber!80!black, font=\scriptsize\bfseries] at (5.15,3.99)
    {Волна II --- наблюдение:\\ наблюдатель внутри};
  \node[anchor=west] at (5.15,3.02) {1970-е фон Фёрстер: второй порядок};
  \node[anchor=west] at (5.15,2.74) {1980 Матурана--Варела: аутопоэзис};
  \node[anchor=west, plosgray] at (7.35,1.30) {1986 backprop-сети};
  \node[anchor=west, plosgray] at (7.35,1.00) {1990 нейроморфные};
  \fill[green!12, rounded corners=3pt] (10.9,0.55) rectangle (14.5,4.05);
  \node[anchor=north west, align=left, green!40!black, font=\scriptsize\bfseries] at (11.05,3.99)
    {Волна III --- когерентность:\\ форма выведена};
  \node[anchor=west] at (11.05,2.98) {УГМ: $N{=}7$, Фано, $G_2$};
  \node[anchor=west] at (11.05,2.68) {КК: меры, динамика};
  \node[anchor=west] at (11.05,2.38) {\SYNARC{}: разум};
  \node[anchor=west, font=\scriptsize\bfseries] at (11.05,2.02) {\TALOS{}: тело};
  \node[anchor=west, plosgray] at (11.05,1.66) {(этот документ)};
\end{tikzpicture}
\caption{Три волны кибернетики. Волна I построила машины, держащие курс
(обратная связь, информация, гомеостат), но угадывала каждую структуру.
Волна II поместила наблюдателя внутрь системы (второй порядок, аутопоэзис),
но могла описать само-наблюдающую петлю, не вывести её форму. Волна III
выводит форму --- семь мод, проводка Фано, одна динамика --- и \TALOS{} ---
её железо.}
\label{fig:waves}
\end{figure}

\begin{zerobox}[С нуля --- что такое машина и что такое правление]
Начнём с ничего. \emph{Машина} --- всё, что меняет метки по твёрдым
правилам: счёты двигают костяшки, музыкальная шкатулка превращает штифты в
ноты. Большинство машин бегут \emph{с открытыми глазами, но слепые} ---
действуют и никогда не проверяют. Дай машине один лишний талант: пусть она
\emph{смотрит на разницу} между желаемым и случившимся и давит против этой
разницы. Комнатный термостат делает ровно это: холодно --- греть, жарко ---
выключить. Этот талант зовётся \emph{обратной связью}, и машина с ним держит
курс сквозь мир, который толкается. Это Волна~I одним вдохом. Волна~II
спрашивает: а если машина смотрит и на \emph{себя смотрящую}? А Волна~III
задаёт лучший детский вопрос --- \emph{тогда какой формы она обязана быть?}
--- и отвечает доказательством. Всё остальное в документе --- этот ответ,
распакованный медленно.
\end{zerobox}

\subsection{Как читать этот документ с нуля}
\label{subsec:howtoread}

Документ написан дважды, вперемешку. \emph{Голос спецификации} несёт
определения, теоремы, таблицы и числа --- всё, что нужно инженеру, без
смягчений. Три повторяющихся бокса несут \emph{голос входа}, и вместе они
образуют полный путь, не предполагающий никаких знаний:

\begin{itemize}[leftmargin=1.5em,itemsep=1pt]
  \item \textbf{\textcolor{talosamber!85!black}{Миф}} (янтарный) продолжает
        историю Талоса: каждая серия кладёт один архитектурный факт на
        бронзового исполина --- обход, жилу, гвоздь. Помнишь историю ---
        помнишь машину.
  \item \textbf{\textcolor{plosblue!70!black}{С нуля}} (синий) пересобирает
        идею текущего раздела из бытовых предметов --- игр, игрушек,
        термостатов --- без символов и без предпосылок. Терпеливый читатель
        может пройти \emph{только} по синим боксам, от начала до конца, и
        прийти к пониманию, что это за машина и почему её форма вынуждена.
  \item \textbf{\textcolor{green!45!black}{Попробуй сам}} (зелёный) делает
        заявления осязаемыми: пальцевые упражнения на фигуре Фано, игра в
        три вопроса, запускаемые команды против эталонного эмулятора
        (\texttt{talos\_v0\_emulator.py} в папке \texttt{architecture/}) ---
        интерактивный слой документа. Ничто в зелёном боксе не требует
        особого железа.
\end{itemize}

Рекомендуемый первый проход для новичка: прочесть все янтарные и синие боксы
по порядку, сыграть в две зелёные пальцевые игры (\S\ref{sec:fabric},
\S\ref{sec:ecc}), затем вернуться на первую страницу и читать голос
спецификации с историей уже в руках. Рекомендуемый проход для инженера:
читать голос спецификации насквозь; боксы никуда не денутся, когда
определению захочется образа.
